\chapter{Project Proposal}
\ChapterLabel{ProjectProposal}
\index{project proposal}
\index{DevSecOps}
\index{SWOT}

\section{Title}
\begin{center}
\large\textbf{TwinTrack: A Side-by-Side DevSecOps Comparison of GitHub Actions
and Jenkins Delivering a Containerized Task Tracker to AWS}
\end{center}

\section{Project Description}
Group~1 will build one small but complete web application, \emph{TwinTrack},
a team task tracker, and deliver it to \gls{aws} through \textbf{two parallel
DevSecOps pipelines} built with comparable tools: \textbf{GitHub Actions}
(hosted, YAML-defined) and \textbf{Jenkins} (self-hosted, Groovy
\texttt{Jenkinsfile}). Both pipelines start from the same commit in the same
GitHub repository, run the same tests and security scans, build the same Docker
image, and deploy it to identical AWS environments (``blue'' for GitHub Actions,
``green'' for Jenkins). Because the application, tests, and target are held
constant, differences in speed, reliability, security coverage, cost, and
operator effort can be attributed to the pipeline tool itself.

An operations dashboard will collect metrics from both pipelines and both
running environments, and the group will finish with a modified \gls{swot}
chart for each tool and a recommendation for which tool suits which kind of
team.

\subsection{The Application}
TwinTrack is deliberately simple so the effort goes into DevOps rather than
features:
\begin{itemize}
\item A REST API (Node.js and Express, the stack from the two-button site) to
create, list, update, and close tasks, with a \texttt{/health} endpoint and a
Prometheus \texttt{/metrics} endpoint.
\item A small browser front end written with class-based JavaScript, following
the design in \ChapterRef{ClassBasedWebsite}.
\item A PostgreSQL database for tasks and users.
\end{itemize}

\subsection{Sample Tasks}
\begin{tabularx}{\linewidth}{L{1.0cm}YL{3.0cm}}
\caption{Sample Project Tasks\label{Table:ProposalTasks}}\\
\toprule \textbf{ID} & \textbf{Task} & \textbf{Area}\\ \midrule
\endfirsthead
\toprule \textbf{ID} & \textbf{Task} & \textbf{Area}\\ \midrule
\endhead \bottomrule \endlastfoot
T1 & Set up the GitHub repository, branch protection on \texttt{main}, pull-request template, and GitHub Projects Kanban board & Source control\\
T2 & Build the TwinTrack API and front end with Docker Compose for local development (app + PostgreSQL) & Development\\
T3 & Write unit, API, and end-to-end tests; enforce a coverage threshold & Testing\\
T4 & Provision two identical AWS environments (blue and green) with Terraform & Infrastructure as code\\
T5 & Write the GitHub Actions workflow: lint, test, scan, build, push to Amazon ECR, deploy to blue & CI/CD tool 1\\
T6 & Run Jenkins in Docker on EC2 and write an equivalent \texttt{Jenkinsfile} that deploys to green & CI/CD tool 2\\
T7 & Add security gates to both pipelines: secret scanning, dependency audit, container image scanning, and a DAST baseline scan & Security\\
T8 & Build the operations dashboard in Grafana (pipeline and runtime metrics) and compare with Amazon CloudWatch & Monitoring\\
T9 & Run the planted-fault experiment (Section~\ref{Section:Unique}) against both pipelines & Experiment\\
T10 & Write SWOT charts, comparison tables, and the final report; build the report with the LaTeX container from \ChapterRef{LatexDocker} & Documentation\\
\end{tabularx}

\section{DevSecOps Tools}
At this stage the list is intentionally broad; tools may be swapped as the
group learns more. Pairs marked ``vs.'' are candidates for head-to-head
comparison.

\begin{tabularx}{\linewidth}{L{2.8cm}YL{3.0cm}}
\caption{Planned DevSecOps Tool Chain\label{Table:ProposalTools}}\\
\toprule \textbf{Category} & \textbf{Tool(s)} & \textbf{Role}\\ \midrule
\endfirsthead
\toprule \textbf{Category} & \textbf{Tool(s)} & \textbf{Role}\\ \midrule
\endhead \bottomrule \endlastfoot
Source control & Git, GitHub (branch protection, pull requests, CODEOWNERS) & Single source of truth for both pipelines\\
Planning & GitHub Projects Kanban board & Track tasks T1--T10\\
CI/CD (\textbf{primary comparison}) & \textbf{GitHub Actions vs.\ Jenkins} & Build, test, scan, deploy\\
Containers & Docker, Docker Compose & Package and run app and database\\
Image registry & Amazon ECR & Store versioned images\\
Infrastructure as code & Terraform & Create identical blue and green environments\\
Hosting & AWS EC2 (Docker on Amazon Linux 2023); possible stretch to ECS Fargate & Run the application\\
Database & PostgreSQL (container locally, Amazon RDS in AWS) & Persist tasks\\
Unit and API testing & Node \texttt{node:test} or Jest, Supertest & Fast feedback on every commit\\
End-to-end testing & Playwright & Test the UI in a real browser\\
Load testing & k6 & Measure latency under load\\
Secret scanning & Gitleaks & Block committed secrets\\
Dependency scanning & \texttt{npm audit}, GitHub Dependabot & Find vulnerable packages\\
Image scanning & Trivy & Find OS and library CVEs in images\\
Static analysis & ESLint, GitHub CodeQL & Code quality and security bugs\\
Dynamic analysis & OWASP ZAP baseline scan & Test the running site\\
Secrets management & GitHub Encrypted Secrets vs.\ Jenkins Credentials; AWS IAM roles & Keep keys out of code\\
Monitoring (\textbf{secondary comparison}) & Prometheus + Grafana vs.\ Amazon CloudWatch & Runtime and pipeline metrics\\
Documentation & LaTeX in Docker (\ChapterRef{LatexDocker}) & Reproducible report builds\\
\end{tabularx}

\section{Pipeline Design}
Both pipelines implement the same stages so they can be compared stage by stage:
\begin{minted}{text}
commit/PR -> checkout -> secret scan (Gitleaks) -> install -> lint (ESLint)
  -> unit + API tests -> dependency audit -> docker build -> image scan (Trivy)
  -> push to ECR -> deploy to blue|green -> smoke test (/health)
  -> DAST (ZAP baseline) -> publish metrics
\end{minted}
A pipeline fails if any security gate reports a high or critical finding.
Pipeline timing, result, and finding counts are exported to the dashboard.

\section{Environment}
\begin{itemize}
\item \textbf{Cloud:} AWS, region \texttt{us-east-1}, using the account,
security model, and scripts established in \ChapterRef{AWSDockerDeployment}.
\item \textbf{Blue environment:} EC2 instance deployed by GitHub Actions
through an OpenID Connect (OIDC) role, so no long-lived AWS keys are stored in
GitHub.
\item \textbf{Green environment:} EC2 instance deployed by Jenkins, which runs
in its own Docker container on a separate EC2 instance with an IAM instance
role.
\item \textbf{Local:} Docker Desktop and Docker Compose on members' laptops.
\item \textbf{Budget:} Free Tier instances where possible, an AWS Budget alert,
and Terraform \texttt{destroy} after each work session.
\end{itemize}

\section{Operations Dashboard}
\index{dashboard}
A Grafana dashboard will place both pipelines and both environments side by
side. Planned panels:
\begin{itemize}
\item \textbf{\gls{dora} metrics per tool:} deployment frequency, lead time for
changes, change failure rate, and time to restore service.
\item \textbf{Pipeline health:} median and 95th-percentile duration per stage,
queue time, and success rate.
\item \textbf{Security:} open findings by scanner and severity, and time from
introducing a planted fault to the pipeline blocking it.
\item \textbf{Runtime:} request rate, error rate, and latency from the app's
\texttt{/metrics} endpoint; CPU and memory of each instance.
\item \textbf{Cost:} estimated cost per deployment and monthly cost per
environment (from AWS Cost Explorer).
\end{itemize}
The same signals will be reproduced in a CloudWatch dashboard to compare the
managed AWS option with the self-hosted Prometheus and Grafana option.

\section{Comparison Criteria}
\begin{tabularx}{\linewidth}{L{2.6cm}Y}
\caption{How the Two CI/CD Tools Will Be Compared\label{Table:ComparisonCriteria}}\\
\toprule \textbf{Dimension} & \textbf{Measure}\\ \midrule
\endfirsthead
\toprule \textbf{Dimension} & \textbf{Measure}\\ \midrule
\endhead \bottomrule \endlastfoot
Security & Planted faults caught; secret-handling model; attack surface of the CI server; supply-chain risk of plugins or marketplace actions\\
Development & Lines of pipeline code; time to write and debug the pipeline; local testability\\
Hosting & Who runs the CI server; setup time; upgrade effort\\
Monitoring & Metrics available out of the box; effort to export to Grafana\\
Testing & Ease of parallel test jobs and test-report publishing\\
Operations & Pipeline duration; failure and recovery time; maintenance hours per week\\
Cost & AWS and CI minutes per deployment; cost at the group's scale and projected at a larger team's scale\\
\end{tabularx}

\section{Preliminary Modified SWOT Charts}
The ``modified'' SWOT used here keeps the four classic quadrants but scores
each one from the perspective of a small DevSecOps team and adds a final
\emph{Evidence to collect} row, so every claim is tested during the project
rather than assumed. These are initial hypotheses and will be revised with
measured results in the final report. The four charts are
Tables~\ref{Table:SWOT-GitHub Actions}--\ref{Table:SWOT-Amazon CloudWatch}.

\newcommand{\SWOTChart}[6]{%
\begin{table}[htbp]
\centering
\caption{Preliminary Modified SWOT: #1}\label{Table:SWOT-#1}
\begin{tabular}{|p{0.455\linewidth}|p{0.455\linewidth}|}
\hline
\textbf{Strengths}\par #2 & \textbf{Weaknesses}\par #3\\ \hline
\textbf{Opportunities}\par #4 & \textbf{Threats}\par #5\\ \hline
\multicolumn{2}{|p{\dimexpr0.91\linewidth+2\tabcolsep+\arrayrulewidth\relax}|}{\textbf{Evidence to collect:} #6}\\ \hline
\end{tabular}
\end{table}}

\SWOTChart{GitHub Actions}
{No server to run; tight GitHub integration; large marketplace; OIDC to AWS without stored keys.}
{Vendor lock-in to GitHub; limited free minutes for private repos; harder to debug locally.}
{Reusable workflows across courses; CodeQL and Dependabot built in; matrix builds.}
{Compromised third-party actions (supply chain); outages outside the team's control.}
{Pipeline duration, minutes used, planted-fault detection time, lines of YAML.}

\SWOTChart{Jenkins}
{Full control and no usage limits; mature plugin ecosystem; runs anywhere, including air-gapped networks.}
{Team must host, patch, and secure the server; plugin version conflicts; Groovy learning curve.}
{Deep customization; shared libraries; on-premises compliance scenarios.}
{Unpatched server or plugins as an attack surface; credential leakage from misconfigured jobs.}
{Setup and maintenance hours, server cost, CVEs in installed plugins, pipeline duration.}

\SWOTChart{Prometheus and Grafana}
{Open source; powerful queries; one dashboard for both pipelines and both environments.}
{Another service to host and secure; storage retention must be managed.}
{Custom DORA panels; alerting to Slack or e-mail.}
{Dashboard exposed publicly without authentication; metric gaps if exporters fail.}
{Setup time, resource use, number of signals shown without custom code.}

\SWOTChart{Amazon CloudWatch}
{Managed; EC2 metrics with no setup; integrated alarms.}
{Per-metric and per-dashboard cost; weak view of CI pipeline data; AWS-only.}
{Logs Insights for incident analysis; native alarms for auto-recovery.}
{Costs grow with custom metrics; lock-in to AWS.}
{Monthly cost, setup time, ability to show pipeline metrics.}

\section{What Makes This Project Unique}
\label{Section:Unique}
\begin{enumerate}
\item \textbf{Controlled ``twin'' experiment.} The same commit, tests, image,
and AWS target run through both tools at the same time, so the comparison is
measured rather than anecdotal.
\item \textbf{Planted-fault security drill.} On a test branch, the group will
deliberately introduce known problems one at a time: a fake (non-functional)
AWS key in source, a dependency with a published CVE, a base image with known
vulnerabilities, and a missing security header. The dashboard records which
pipeline catches each fault, at which stage, and how long it takes. This turns
``which tool is more secure?'' into a scored result.
\item \textbf{Cost per deployment.} Each successful deployment is tagged with
its CI time and AWS resources, giving a cost-per-deploy figure for each tool.
\item \textbf{Docs as code.} The final report is built from the repository by
the LaTeX container, so the report itself can be a pipeline artifact.
\end{enumerate}

\section{Mapping to the Project Rubric}
\begin{tabularx}{\linewidth}{L{5.0cm}Y}
\caption{Proposal Coverage of the Project Expectations\label{Table:RubricMapping}}\\
\toprule \textbf{Expectation} & \textbf{How TwinTrack addresses it}\\ \midrule
\endfirsthead
\toprule \textbf{Expectation} & \textbf{How TwinTrack addresses it}\\ \midrule
\endhead \bottomrule \endlastfoot
Selected tools and environment & Table~\ref{Table:ProposalTools}; AWS plus local Docker\\
Sample working application & TwinTrack API, front end, and PostgreSQL database\\
Specific DevOps tools to build it & Git/GitHub, Docker, Terraform, Trivy, Gitleaks, ZAP, Playwright, k6\\
Two comparable tools & GitHub Actions vs.\ Jenkins (primary); Prometheus/Grafana vs.\ CloudWatch (secondary)\\
Comparison across security, development, hosting, monitoring, testing, operations & Table~\ref{Table:ComparisonCriteria}\\
Environment (AWS, Digital Ocean, or homebrew) & AWS EC2, ECR, RDS; Docker locally\\
Operational dashboard & Grafana dashboard with DORA, pipeline, security, runtime, and cost panels\\
Modified SWOT chart for each tool & Preliminary charts above; final charts with measured evidence\\
Something unique & Twin experiment, planted-fault drill, cost per deploy, docs as code\\
\end{tabularx}

\section{Team Roles and Timeline}
\begin{tabularx}{\linewidth}{L{2.6cm}Y}
\toprule \textbf{Member} & \textbf{Primary responsibility}\\ \midrule
Brahmi Dalvi & Application, tests, GitHub Actions pipeline, Grafana dashboard\\
Sakshi Deshpande & Terraform/AWS environments, Jenkins pipeline, security gates, report build\\
Both & Planted-fault experiment, SWOT analysis, final report and presentation\\ \bottomrule
\end{tabularx}

\medskip
\begin{tabularx}{\linewidth}{L{2.6cm}Y}
\toprule \textbf{Weeks} & \textbf{Milestone}\\ \midrule
1--2 & Repository, Kanban board, app skeleton, Docker Compose (T1--T2)\\
3--4 & Tests and Terraform environments (T3--T4)\\
5--6 & GitHub Actions and Jenkins pipelines with security gates (T5--T7)\\
7 & Dashboard and planted-fault experiment (T8--T9)\\
8 & Final SWOT charts, report, and presentation (T10)\\ \bottomrule
\end{tabularx}
